% !PRINTED 86
\weberpage{86}

%<<CONT
in which the functions $f(\omega)$, $f_1(\omega)$, $f_2(\omega)$
are defined as follows:
\weq{11}{
\begin{aligned}
f(\omega) &= q^{-\frac{1}{24}}\prod_{1,\infty}^{\nu}(1+q^{2\nu-1}),\\
f_1(\omega) &= q^{-\frac{1}{24}}\prod_{1,\infty}^{\nu}(1-q^{2\nu-1}),\\
f_2(\omega) &= \sqrt{2}\,q^{\frac{1}{12}}\prod_{1,\infty}^{\nu}(1+q^{2\nu}).
\end{aligned}
}

The functions $\eta(\omega)$, $f(\omega)$, $f_1(\omega)$,
$f_2(\omega)$ will play an important part in our later
investigations. From § 21, (14) there follows, by (10),
the relation
\weq{12}{f(\omega)^8 = f_1(\omega)^8 + f_2(\omega)^8,}
and from § 23 (5), by (9),
\weq{13}{f(\omega)f_1(\omega)f_2(\omega) = \sqrt{2}.}
The last formula also follows from (11) by the use of the
identical relation:
\[ \prod (1+q^\nu)(1-q^{2\nu-1}) = \prod \frac{(1+q^\nu)(1-q^\nu)}{1-q^{2\nu}} = 1. \]

\begin{center}
\textbf{§ 25. Representation of the $\vartheta$-Functions by
Infinite Series.}
\end{center}

The second way of arriving at the representation of the
$\vartheta$-functions consists in seeking to form a convergent
infinite series that satisfies the fundamental equations.

Let us first observe that, on account of the condition
\[ \vartheta_{00}(v+1) = \vartheta_{00}(v) \]
the function $\vartheta_{00}(v)$ can be regarded as a one-valued
function of $e^{2\pi i v}$, and let us accordingly put
\[ \vartheta_{00}(v) = \sum_{-\infty,\infty}^{\nu} A_\nu e^{2\pi i v \nu}, \]
the differential equation § 20, (4):
\[ \frac{\partial^2 \vartheta}{\partial v^2} - 4\pi i \frac{\partial \vartheta}{\partial \omega} = 0 \]
then gives for $A_\nu$ the condition
\[ \frac{dA_\nu}{d\omega} = \pi i \nu^2 A_\nu, \qquad A_\nu = c_\nu e^{\pi i \omega \nu^2} = c_\nu q^{\nu^2}, \]
in which $c_\nu$ is constant with respect to $\omega$. Thus we obtain
\[ \vartheta_{00}(v) = \sum c_\nu q^{\nu^2} c^{2\pi i v \nu}, \]
% [sic] German source prints c as the exponential base; mathematically expected e.
%CONT>>
